\chapter{2007年大纲II卷（理）}
\section{单选题}
\begin{problem}
\(\sin 210^{\circ} =\)（\quad）

\choices
    {\(\frac{\sqrt{3}}{2}\)}
    {\(-\frac{\sqrt{3}}{2}\)}
    {\( \frac{1}{2} \)}
    {\(-\frac{1}{2}\)}
\end{problem}

\begin{answer}
D
\end{answer}

\begin{solution}
由诱导公式，
\[
\sin 210^{\circ}=\sin(180^{\circ}+30^{\circ})=-\sin30^{\circ}=-\frac{1}{2}
\]
因此选 D.
\end{solution}

\begin{problem}
函数 \( f(x) = |\sin x| \) 的一个单调递增区间是（\quad）

\choices
    {\(\left(-\frac{\pi}{4}, \frac{\pi}{4}\right)\)}
    {\(\left(\frac{\pi}{4}, \frac{3\pi}{4}\right)\)}
    {\(\left(\pi, \frac{3\pi}{2}\right)\)}
    {\(\left(\frac{3\pi}{2}, 2\pi\right)\)}
\end{problem}

\begin{answer}
C
\end{answer}

\begin{solution}
在 \(\left(-\frac{\pi}{4},0\right)\) 内，\(f(x)=-\sin x\)，且 \(f'(x)=-\cos x<0\)；在 \(\left(0,\frac{\pi}{4}\right)\) 内，\(f(x)=\sin x\)，且 \(f'(x)=\cos x>0\)，所以选项 A 中的函数先减后增. 在 \(\left(\frac{\pi}{4},\frac{3\pi}{4}\right)\) 内，\(f(x)=\sin x\)，而 \(\cos x\) 在 \(\frac{\pi}{2}\) 两侧变号，函数先增后减. 在 \(\left(\pi,\frac{3\pi}{2}\right)\) 内，\(f(x)=-\sin x\)，且 \(-\cos x>0\)，故函数单调递增. 在 \(\left(\frac{3\pi}{2},2\pi\right)\) 内，\(-\cos x<0\)，故函数单调递减. 因此选 C.
\end{solution}

\begin{problem}
设复数 \(z\) 满足 \(\frac{1 + 2\i}{z} = \i\)，则 \(z =\)（\quad）

\choices
    {\(-2 + \i\)}
    {\(-2 - \i\)}
    {\(2 - \i\)}
    {\(2 + \i\)}
\end{problem}

\begin{answer}
C
\end{answer}

\begin{solution}
因分母不为零，且 \(\frac{1+2\i}{z}=\i\)，所以 \(z=\frac{1+2\i}{\i}=(1+2\i)(-\i)=2-\i\). 因此选 C.
\end{solution}

\begin{problem}
以下四个数中的最大者是（\quad）

\choices
    {\((\ln 2)^{2}\)}
    {\(\ln (\ln 2)\)}
    {\(\ln \sqrt{2}\)}
    {\(\ln 2\)}
\end{problem}

\begin{answer}
D
\end{answer}

\begin{solution}
令 \(u=\ln2\). 因为 \(1<2<\e\)，所以 \(0<u<1\). 从而 \(u^2<u\)，\(\ln u<0\)，且 \(\ln\sqrt2=\frac12\ln2=\frac{u}{2}<u\). 因此四个数中最大的是 \(\ln2\)，即选 D.
\end{solution}

\begin{problem}
在 \(\triangle ABC\) 中，已知 \(D\) 是 \(AB\) 边上一点，若 \(\overrightarrow{AD} = 2\overrightarrow{DB}\)，\(\overrightarrow{CD} = \frac{1}{3}\overrightarrow{CA} + \lambda \overrightarrow{CB}\)，则 \(\lambda =\)（\quad）

\choices
    {\( \frac{2}{3} \)}
    {\(\frac{1}{3}\)}
    {\(-\frac{1}{3}\)}
    {\(-\frac{2}{3}\)}
\end{problem}

\begin{answer}
A
\end{answer}

\begin{solution}
由 \(\overrightarrow{AD}=2\overrightarrow{DB}\)，得 \(\overrightarrow{AB}=\overrightarrow{AD}+\overrightarrow{DB}=3\overrightarrow{DB}\)，所以 \(\overrightarrow{AD}=\frac{2}{3}\overrightarrow{AB}\). 又 \(\overrightarrow{AB}=\overrightarrow{AC}+\overrightarrow{CB}=-\overrightarrow{CA}+\overrightarrow{CB}\)，故
\[
\overrightarrow{CD}=\overrightarrow{CA}+\overrightarrow{AD}
=\overrightarrow{CA}+\frac{2}{3}\left(-\overrightarrow{CA}+\overrightarrow{CB}\right)
=\frac{1}{3}\overrightarrow{CA}+\frac{2}{3}\overrightarrow{CB}
\]
故 \(\lambda=\frac{2}{3}\)，选 A.
\end{solution}

\begin{problem}
不等式 \(\frac{x - 1}{x^2 - 4} > 0\) 的解集为（\quad）

\choices
    {\((-2, 1)\)}
    {\((2, +\infty)\)}
    {\((-2,1)\cup (2, + \infty)\)}
    {\((- \infty, -2) \cup (1, +\infty)\)}
\end{problem}

\begin{answer}
C
\end{answer}

\begin{solution}
分式有意义的条件为 \(x\ne-2\)，\(2\). 分子、分母的零点依次为 \(1\) 和 \(-2,2\). 按 \(-2<1<2\) 分段：在 \((-\infty,-2)\) 和 \((1,2)\) 内，分子、分母异号，分式为负；在 \((-2,1)\) 和 \((2,+\infty)\) 内，分子、分母同号，分式为正. 因而
\[
\frac{x-1}{x^2-4}>0\quad\Longleftrightarrow\quad -2<x<1\ \text{或}\ x>2
\]
严格不等式不取 \(x=1\)，分母零点也不能取，所以解集为 \((-2,1)\cup(2,+\infty)\)，选 C.
\end{solution}

\begin{problem}
已知正三棱柱 \(ABC - A_{1}B_{1}C_{1}\) 的侧棱长与底面边长相等，则 \(AB_{1}\) 与侧面 \(ACC_{1}A_{1}\) 所成角的正弦等于（\quad）

\choices
    {\(\frac{\sqrt{6}}{4}\)}
    {\(\frac{\sqrt{10}}{4}\)}
    {\(\frac{\sqrt{2}}{2}\)}
    {\(\frac{\sqrt{3}}{2}\)}
\end{problem}

\begin{answer}
A
\end{answer}

\begin{solution}
设正三棱柱的底面边长和侧棱长均为 \(a>0\). 建立空间直角坐标系，使
\(A=(0,0,0)\)，\(C=(a,0,0)\)，\(B=\left(\frac a2,\frac{\sqrt3 a}{2},0\right)\)
且 \(A_1\)，\(B_1\)，\(C_1\) 的竖直坐标均为 \(a\). 于是
\(\overrightarrow{AB_1}=\left(\frac a2,\frac{\sqrt3 a}{2},a\right)\)，\(|AB_1|=\sqrt2a\)
侧面 \(ACC_1A_1\) 的方程为 \(y=0\)，其法向量可取 \(\bs{n}=(0,1,0)\). 若直线 \(AB_1\) 与该平面所成角为 \(\theta\)，则
\[
\sin\theta=\frac{|\overrightarrow{AB_1}\cdot\bs{n}|}{|\overrightarrow{AB_1}|\,|\bs{n}|}
=\frac{\frac{\sqrt3a}{2}}{\sqrt2a}=\frac{\sqrt6}{4}
\]
因此选 A.
\end{solution}

\begin{problem}
已知曲线 \( y = \frac{x^2}{4} - 3\ln x \) 的一条切线的斜率为 \( \frac{1}{2} \)，则切点的横坐标为（\quad）

\choices
    {\(3\)}
    {\(2\)}
    {\(1\)}
    {\(\frac{1}{2}\)}
\end{problem}

\begin{answer}
A
\end{answer}

\begin{solution}
函数定义域为 \(x>0\)，且
\[
y'=\frac{x}{2}-\frac{3}{x}
\]
由切线斜率为 \(\frac12\)，得
\[
\frac{x}{2}-\frac{3}{x}=\frac12
\quad\Longleftrightarrow\quad x^2-x-6=0
\quad\Longleftrightarrow\quad (x-3)(x+2)=0
\]
结合 \(x>0\)，切点横坐标为 \(3\)，选 A.
\end{solution}

\begin{problem}
把函数 \( y = \e^{x} \) 的图象按向量 \( \bs{a} = (2,3) \) 平移，得到 \( y = f(x) \) 的图象，则 \( f(x) = \)（\quad）

\choices
    {\(\e^{x - 3} + 2\)}
    {\(\e^{x + 3} - 2\)}
    {\(\e^{x - 2} + 3\)}
    {\(\e^{x + 2} - 3\)}
\end{problem}

\begin{answer}
C
\end{answer}

\begin{solution}
原图象上的点可写为 \((u,\e^u)\). 平移向量 \((2,3)\) 后变为 \((x,y)=(u+2,\e^u+3)\)，即 \(u=x-2\). 因此平移后的函数为
\[
y=\e^{x-2}+3
\]
故 \(f(x)=\e^{x-2}+3\)，选 C.
\end{solution}

\begin{problem}
从 \(5\) 位同学中选派 \(4\) 位同学在星期五，星期六，星期日参加公益活动，每人一天，要求星期五有 \(2\) 人参加，星期六，星期日各有 \(1\) 人参加，则不同的选派方法共有（\quad）

\choices
    {\(40\) 种}
    {\(60\) 种}
    {\(100\) 种}
    {\(120\) 种}
\end{problem}

\begin{answer}
B
\end{answer}

\begin{solution}
先从 \(5\) 位同学中选出星期五参加活动的 \(2\) 人，有 \(\mathrm C_5^2\) 种方法. 再从剩下的 \(3\) 人中选出星期六、星期日参加活动的同学，并区分两天，有 \(\mathrm A_3^2\) 种方法. 因此总数为
\[
\mathrm C_5^2\mathrm A_3^2=10\times6=60
\]
故选 B.
\end{solution}

\begin{problem}
设 \(F_{1}\)，\(F_{2}\) 分别是双曲线 \(\frac{x^{2}}{a^{2}} - \frac{y^{2}}{b^{2}} = 1\) 的左，右焦点. 若双曲线上存在点 \(A\)，使 \(\angle F_{1}AF_{2} = 90^{\circ}\)，且 \(|AF_{1}| = 3|AF_{2}|\)，则双曲线离心率为（\quad）

\choices
    {\(\frac{\sqrt{5}}{2}\)}
    {\(\frac{\sqrt{10}}{2}\)}
    {\(\frac{\sqrt{15}}{2}\)}
    {\(\sqrt{5}\)}
\end{problem}

\begin{answer}
B
\end{answer}

\begin{solution}
双曲线的焦距为 \(F_1F_2=2c\). 因为 \(|AF_1|>|AF_2|\)，点 \(A\) 不在左支，而在右支；于是由双曲线定义
\[
|AF_1|-|AF_2|=2a
\]
故 \(|AF_2|=a\)，\(|AF_1|=3a\). 又 \(\angle F_1AF_2=90^{\circ}\)，在直角三角形 \(F_1AF_2\) 中
\[
(2c)^2=(3a)^2+a^2=10a^2
\]
所以离心率
\[
e=\frac ca=\frac{\sqrt{10}}{2}
\]
故选 B.
\end{solution}

\begin{problem}
设 \(F\) 为抛物线 \(y^{2} = 4x\) 的焦点，\(A\)，\(B\)，\(C\) 为该抛物线上三点. 若 \(\overrightarrow{FA} + \overrightarrow{FB} + \overrightarrow{FC} = 0\)，则 \(|FA| + |FB| + |FC| =\)（\quad）

\choices
    {\(9\)}
    {\(6\)}
    {\(4\)}
    {\(3\)}
\end{problem}

\begin{answer}
B
\end{answer}

\begin{solution}
抛物线 \(y^2=4x\) 的焦点为 \(F=(1,0)\). 设抛物线上三点分别对应参数 \(u\)，\(v\)，\(w\)，则它们的坐标为 \((u^2,2u)\)、\((v^2,2v)\)、\((w^2,2w)\)，从而
\(\overrightarrow{FA}=(u^2-1,2u)\)，\(\overrightarrow{FB}=(v^2-1,2v)\)，\(\overrightarrow{FC}=(w^2-1,2w)\)
由向量和为零，得
\(u^2+v^2+w^2=3\)，\(u+v+w=0\)
并且
\[
|FA|=\sqrt{(u^2-1)^2+4u^2}=u^2+1
\]
同理 \(|FB|=v^2+1\)，\(|FC|=w^2+1\). 因此
\[
|FA|+|FB|+|FC|=u^2+v^2+w^2+3=6
\]
故选 B.
\end{solution}

\section{填空题}
\begin{problem}
\((1 + 2x^{2})\left(x - \frac{1}{x}\right)^{8}\) 的展开式中常数项为\fillinblank{}.
\end{problem}

\begin{answer}
\(-42\)
\end{answer}

\begin{solution}
在 \(\left(x-\frac1x\right)^8\) 的展开式中，通项为
\[
(-1)^k\mathrm C_8^k x^{8-2k}
\]
与 \(1\) 相乘得到常数项的条件为 \(8-2k=0\)，此时贡献为 \(\mathrm C_8^4=70\)；与 \(2x^2\) 相乘得到常数项的条件为 \(8-2k+2=0\)，此时贡献为 \(2(-1)^5\mathrm C_8^5=-112\). 所以常数项为
\[
70-112=-42
\]
\end{solution}

\begin{problem}
在某项测量中，测量结果 \(\xi\) 服从正态分布 \(N(1, \sigma^2)\)（\(\sigma > 0\)）. 若 \(\xi\) 在 \((0, 1)\) 内取值的概率为 \(0.4\)，则 \(\xi\) 在 \((0, 2)\) 内取值的概率为\fillinblank{}.
\end{problem}

\begin{answer}
\(0.8\)
\end{answer}

\begin{solution}
正态分布 \(N(1,\sigma^2)\) 的密度图象关于直线 \(x=1\) 对称，所以
\[
P(1<\xi<2)=P(0<\xi<1)=0.4
\]
由于连续型随机变量在单点处的概率为零，故
\[
P(0<\xi<2)=P(0<\xi<1)+P(1<\xi<2)=0.4+0.4=0.8
\]
\end{solution}

\begin{problem}
一个正四棱柱的各个顶点在一个直径为 \(2 \mathrm{~cm}\) 的球面上. 如果正四棱柱的底面边长为 \(1 \mathrm{~cm}\)，那么该棱柱的表面积为\fillinblank{}\(\mathrm{cm}^{2}\).
\end{problem}

\begin{answer}
\(2 + 4\sqrt{2}\)
\end{answer}

\begin{solution}
设正四棱柱的高为 \(h\). 正四棱柱的外接球球心是棱柱中心，因此任意一对相对顶点连成的空间对角线是球的直径. 其底面边长为 \(1\)，所以空间对角线长为 \(\sqrt{2+h^2}\)，从而
\(\sqrt{2+h^2}=2\)，\(h=\sqrt2\)
因此棱柱表面积为
\[
2\times1^2+4\times1\times\sqrt2=2+4\sqrt2\ \mathrm{cm}^2
\]
\end{solution}

\begin{problem}
已知数列的通项 \(a_{n} = -5n + 2\)，其前 \(n\) 项和为 \(S_{n}\)，则 \(\displaystyle\lim_{n\to\infty}\frac{S_n}{n^2} = \)\fillinblank{}.
\end{problem}

\begin{answer}
\(-\frac{5}{2}\)
\end{answer}

\begin{solution}
由等差数列求和公式，
\[
S_n=\sum_{k=1}^n(-5k+2)=-5\frac{n(n+1)}{2}+2n=-\frac52n^2-\frac12n
\]
所以
\[
\lim_{n\to\infty}\frac{S_n}{n^2}
=\lim_{n\to\infty}\left(-\frac52-\frac{1}{2n}\right)=-\frac52
\]
\end{solution}

\section{解答题}
\begin{problem}
在 \(\triangle ABC\) 中，已知内角 \(A = \frac{\pi}{3}\)，边 \(BC = 2\sqrt{3}\)，设内角 \(B = x\)，周长为 \(y\).
\begin{enumerate}
\item 求函数 \( y = f(x) \) 的解析式和定义域；
\item 求 \(y\) 的最大值.
\end{enumerate}
\end{problem}

\begin{solution}
\begin{enumerate}
\item 设 \(a=BC=2\sqrt3\)，由正弦定理，
\(\frac{BC}{\sin A}=\frac{CA}{\sin x}=\frac{AB}{\sin C}\)，\(C=\frac{2\pi}{3}-x\)
由于 \(\sin A=\frac{\sqrt3}{2}\)，得
\(CA=4\sin x\)，\(AB=4\sin\left(\frac{2\pi}{3}-x\right)\)
故周长函数为
\[
y=f(x)=2\sqrt3+4\sin x+4\sin\left(\frac{2\pi}{3}-x\right)
\]
三角形内角必须为正，所以 \(0<x<\frac{2\pi}{3}\).

\item 利用和差化积公式，
\[
\sin x+\sin\left(\frac{2\pi}{3}-x\right)=2\sin\frac{\pi}{3}\cos\left(x-\frac{\pi}{3}\right)=\sqrt3\cos\left(x-\frac{\pi}{3}\right)
\]
因此
\[
y=2\sqrt3+4\sqrt3\cos\left(x-\frac{\pi}{3}\right)\leqslant6\sqrt3
\]
当 \(x=\frac\pi3\) 时等号成立，故 \(y\) 的最大值为 \(6\sqrt3\).
\end{enumerate}
\end{solution}

\begin{problem}
从某批产品中，有放回地抽取产品 \(2\) 次，每次随机抽取 \(1\) 件，假设事件 \(A\)：“取出的 \(2\) 件产品中至多有 \(1\) 件是二等品”的概率 \(P(A)=0.96\).
\begin{enumerate}
\item 求从该批产品中任取 \(1\) 件是二等品的概率 \(p\)；
\item 若该批产品共有 \(100\) 件，从中任意抽取 \(2\) 件，\(\xi\) 表示取出的 \(2\) 件产品中二等品的件数，求 \(\xi\) 的分布列.
\end{enumerate}
\end{problem}

\begin{solution}
\begin{enumerate}
\item 设任取一件产品为二等品的概率为 \(p\). 事件 \(A\) 的对立事件是两次取出的产品都是二等品. 由于有放回抽取，两次抽取相互独立，所以
\[
1-p^2=P(A)=0.96
\]
由 \(0\leqslant p\leqslant1\)，得 \(p=0.2\).

\item 由 \(p=0.2\)，这 \(100\) 件产品中有 \(20\) 件二等品、\(80\) 件非二等品. 任取 \(2\) 件的总数为 \(\mathrm C_{100}^{2}\)，随机变量 \(\xi\) 的取值为 \(0\)，\(1\)，\(2\). 因此
\[
P(\xi=0)=\frac{\mathrm C_{80}^{2}}{\mathrm C_{100}^{2}}=\frac{316}{495}
\]
\(P(\xi=1)=\frac{\mathrm C_{20}^{1}\mathrm C_{80}^{1}}{\mathrm C_{100}^{2}}=\frac{160}{495}\)，\(P(\xi=2)=\frac{\mathrm C_{20}^{2}}{\mathrm C_{100}^{2}}=\frac{19}{495}\)
三项概率之和为 \(\frac{316+160+19}{495}=1\)，故分布列完整.
\end{enumerate}
\end{solution}

\begin{problem}
如图，在四棱锥 \(S - ABCD\) 中，底面 \(ABCD\) 为正方形，侧棱 \(SD \perp\) 底面 \(ABCD\)，\(E\)，\(F\) 分别是 \(AB\)，\(SC\) 的中点.
\begin{enumerate}
\item 求证： \(EF\parallel\) 平面 \(SAD\)；
\item 设 \(SD = 2CD\)，求二面角 \(A - EF - D\) 的大小.
\end{enumerate}
\bitmapfigure[width=0.34\linewidth]{img/2007/syllabus_paper_2_science/q19_fig1.png}
\end{problem}

\begin{solution}
\begin{enumerate}
\item 设正方形边长为 \(a>0\)，建立空间直角坐标系，令
\(D=(0,0,0)\)，\(A=(a,0,0)\)，\(C=(0,a,0)\)，\(B=(a,a,0)\)，\(S=(0,0,h)\)
其中 \(h=SD\). 由中点公式，
\(E=\left(a,\frac a2,0\right)\)，\(F=\left(0,\frac a2,\frac h2\right)\)
于是 \(EF\) 的方向向量为
\[
\overrightarrow{EF}=\left(-a,0,\frac h2\right)
\]
而平面 \(SAD\) 的方程为 \(y=0\). 直线 \(EF\) 的方向向量的第二个分量为 \(0\)，且 \(EF\) 不在平面 \(SAD\) 内，因此 \(EF\parallel\) 平面 \(SAD\).

\item 由 \(SD=2CD\)，得 \(h=2a\)，从而
\(E=(a,\tfrac a2,0)\)，\(F=(0,\tfrac a2,a)\)
由 \(\overrightarrow{EA}=(0,-\frac a2,0)\)，\(\overrightarrow{EF}=(-a,0,a)\)，得 \(\overrightarrow{EA}\times\overrightarrow{EF}\) 与 \((1,0,1)\) 平行；由 \(\overrightarrow{DE}=(a,\frac a2,0)\)，\(\overrightarrow{DF}=(0,\frac a2,a)\)，得 \(\overrightarrow{DE}\times\overrightarrow{DF}\) 与 \((1,-2,1)\) 平行. 因而平面 \(AEF\) 和 \(DEF\) 的法向量分别可取
\(\bs{n}_1=(1,0,1)\)，\(\bs{n}_2=(1,-2,1)\)
故二面角 \(A-EF-D\) 的平面角记为 \(\theta\)，则
\[
\cos\theta=\frac{|\bs{n}_1\cdot\bs{n}_2|}{|\bs{n}_1||\bs{n}_2|}
=\frac{2}{\sqrt2\sqrt6}=\frac1{\sqrt3}
\]
因为二面角取锐角，
\[
\tan\theta=\frac{\sqrt{1-\cos^2\theta}}{\cos\theta}=\sqrt2
\]
所以二面角的大小为 \(\arctan\sqrt2\).
\end{enumerate}
\end{solution}

\begin{problem}
在直角坐标系 \(xOy\) 中，以 \(O\) 为圆心的圆与直线 \(x - \sqrt{3} y = 4\) 相切.
\begin{enumerate}
\item 求圆 \(O\) 的方程；
\item 圆 \(O\) 与 \(x\) 轴相交于 \(A\)，\(B\) 两点，圆内的动点 \(P\) 使 \(\left|PA\right|\)，\(\left|PO\right|\)，\(\left|PB\right|\) 成等比数列，求 \(\overrightarrow{PA} \cdot \overrightarrow{PB}\) 的取值范围.
\end{enumerate}
\end{problem}

\begin{solution}
\begin{enumerate}
\item 原点到直线 \(x-\sqrt3y=4\) 的距离为
\[
R=\frac{4}{\sqrt{1+3}}=2
\]
因此圆 \(O\) 的方程为 \(x^2+y^2=4\).

\item 设 \(P=(x,y)\)，并记 \(r^2=OP^2=x^2+y^2\). 圆与 \(x\) 轴交于 \(A=(-2,0)\)，\(B=(2,0)\). 由 \(|PA|\)、\(|PO|\)、\(|PB|\) 成等比数列，
\[
|PO|^2=|PA||PB|
\]
两边平方并代入坐标，得
\[
r^4=\bigl((x+2)^2+y^2\bigr)\bigl((x-2)^2+y^2\bigr)
=\bigl(r^2+4+4x\bigr)\bigl(r^2+4-4x\bigr)
\]
化简为
\[
r^2=2x^2-2
\]
又 \(y^2=r^2-x^2=x^2-2\geqslant0\)，所以 \(x^2\geqslant2\). 因 \(P\) 在圆内，\(r^2<4\)，故 \(x^2<3\). 另一方面，
\[
\overrightarrow{PA}\cdot\overrightarrow{PB}
=(-2-x,-y)\cdot(2-x,-y)=x^2+y^2-4=r^2-4=2x^2-6
\]
于是
\[
\overrightarrow{PA}\cdot\overrightarrow{PB}\in[-2,0)
\]
当 \(2\leqslant x^2<3\) 时取 \(y^2=x^2-2\)，则 \(P\) 在圆内且满足等比条件，上述范围中的每个值均可取得，所以该范围即为所求取值范围.
\end{enumerate}
\end{solution}

\begin{problem}
设数列 \(\{a_{n}\}\) 的首项 \(a_1 \in (0,1)\)，\(a_n = \frac{3 - a_{n-1}}{2}\)，\(n = 2\)，\(3\)，\(4\)，\(\cdots\).
\begin{enumerate}
\item 求 \( \{a_{n}\} \) 的通项公式；
\item 设 \( b_{n} = a_{n}\sqrt{3 - 2a_{n}} \)，证明 \( b_{n} < b_{n+1} \)，其中 \( n \) 为正整数.
\end{enumerate}
\end{problem}

\begin{solution}
\begin{enumerate}
\item 由递推式，
\[
a_n-1=\frac{3-a_{n-1}}{2}-1=-\frac12(a_{n-1}-1)
\]
因此 \(\{a_n-1\}\) 是首项 \(a_1-1\)、公比 \(-\frac12\) 的等比数列，
\[
a_n=1+(a_1-1)\left(-\frac12\right)^{n-1}=1-(1-a_1)\left(-\frac12\right)^{n-1}
\]

\item 令 \(g(t)=t^2(3-2t)\). 由通项公式，若 \(n\) 为奇数，则 \(0<(1-a_1)2^{-(n-1)}<1\)，故 \(0<a_n<1\)；若 \(n\) 为偶数，则 \(a_n=1+(1-a_1)2^{-(n-1)}\)，故 \(1<a_n<\frac32\). 因而对任意 \(n\) 都有 \(a_n>0\) 且 \(3-2a_n>0\)，所以 \(b_n>0\). 设 \(t=a_n\)，则 \(a_{n+1}=\frac{3-t}{2}\)，从而
\[
b_{n+1}^2-b_n^2
=g\left(\frac{3-t}{2}\right)-g(t)
=\frac{9}{4}t(t-1)^2
\]
由于 \(a_1\ne1\)，由通项公式知任意正整数 \(n\) 都有 \(a_n\ne1\)，故上式严格大于零. 因此 \(b_{n+1}^2>b_n^2\)，结合 \(b_n>0\)，得到 \(b_n<b_{n+1}\).
\end{enumerate}
\end{solution}

\begin{problem}
已知函数 \( f(x) = x^3 - x \).
\begin{enumerate}
\item 求曲线 \( y = f(x) \) 在点 \( M(t, f(t)) \) 处的切线方程；
\item 设 \(a > 0\)，如果过点 \((a, b)\) 可作曲线 \(y = f(x)\) 的三条切线，证明：\(-a < b < f(a)\).
\end{enumerate}
\end{problem}

\begin{solution}
\begin{enumerate}
\item \(f'(x)=3x^2-1\)，所以在点 \(M(t,f(t))\) 处的切线为
\[
y-f(t)=f'(t)(x-t)
\]
代入 \(f(t)=t^3-t\) 和 \(f'(t)=3t^2-1\)，得
\[
y-(t^3-t)=(3t^2-1)(x-t)
\]
即
\[
y=(3t^2-1)x-2t^3
\]

\item 设过点 \((a,b)\) 的切线与曲线的切点横坐标为 \(t\). 由上面的切线方程，\((a,b)\) 在该切线上等价于
\[
b=(3t^2-1)a-2t^3
\]
即
\[
g(t)=2t^3-3at^2+a+b=0
\]
三条切线对应三个不同的实数根. 又
\[
g'(t)=6t(t-a)
\]
因为 \(a>0\)，函数 \(g(t)\) 在 \((-\infty,0)\) 和 \((a,+\infty)\) 上递增，在 \((0,a)\) 上递减. 要有三个不同的实根，极大值 \(g(0)\) 必须为正、极小值 \(g(a)\) 必须为负；否则单调三段至多产生两个不同的交点. 因而
\(g(0)=a+b>0\)，\(g(a)=b-a^3+a=b-f(a)<0\)
于是
\[
-a<b<f(a)
\]
\end{enumerate}
\end{solution}
